% Hardware, software, and the models of the Day 4 lab.
\begin{frame}{Hardware and software}
  \begin{columns}[T]
    \begin{column}{0.45\textwidth}
      \textbf{Hardware for each group}
      \begin{itemize}
        \item XIAOML Kit with the Arduino firmware. You need it in Part D
              only.
        \item USB-C data cable.
        \item Laptop with your Day 2 dataset in \code{Labs/day02/data/}.
      \end{itemize}
    \end{column}
    \begin{column}{0.51\textwidth}
      \textbf{Software}
      \begin{itemize}
        \item Arduino IDE 2 with the esp32 core 3.3.12.
        \item The library Chirale\_TensorFlowLite 2.0.0 of Day 3.
        \item The Python packages of Day 3: \code{tensorflow-cpu} and
              \code{ai-edge-litert}.
        \item No new package and no new library.
      \end{itemize}
    \end{column}
  \end{columns}
  \vskip 0.6em
  \small
  \renewcommand{\arraystretch}{1.15}
  \begin{tabular}{@{}lL{0.80\textwidth}@{}}
    Lab files & \code{Labs/day04/}. Run the commands from this folder. \\
    Notebook  & \code{../../.venv/bin/jupyter lab quantization.ipynb} \\
    Sketch    & \code{sketches/motion\_quant/motion\_quant.ino} \\
  \end{tabular}
\end{frame}

\begin{frame}{Three objects, one method}
  \begingroup\centering
  \begin{tikzpicture}[font=\footnotesize, >={Stealth[length=5pt]},
      b/.style={draw=coolgray, rounded corners=3pt, fill=boxgray,
                minimum width=2.5cm, minimum height=1.0cm, align=center},
      lab/.style={font=\scriptsize, text=coolgray, align=center}]
    \node[b] (a) at (0,0)    {One tensor};
    \node[b] (b) at (3.6,0)  {A small CNN\\for the raw window};
    \node[b, draw=cardinalred] (c) at (7.2,0) {The feature model\\of Day 3};
    \draw[->] (a) -- (b); \draw[->] (b) -- (c);
    \node[lab] at (0,-0.95)   {Part A\\NumPy};
    \node[lab] at (3.6,-0.95) {Parts B and C\\Keras and LiteRT};
    \node[lab] at (7.2,-0.95) {Part D\\the board};
    \node[lab] at (0,0.85)   {393\,600 values};
    \node[lab] at (3.6,0.85) {852 parameters};
    \node[lab] at (7.2,0.85) {1534 parameters};
  \end{tikzpicture}\par\endgroup
  \vskip 0.5em
  \small
  \begin{itemize}
    \item \textbf{Part A:} you write the mapping for one tensor: the
          acceleration values of your training windows.
    \item \textbf{Parts B and C:} the converter and the training apply the
          same mapping to each tensor of a CNN. The CNN takes the raw
          window in $g$, so the range of its input depends on the motion.
    \item \textbf{Part D:} the model of Day 3 goes to the board two times:
          in \code{float32} and in \code{int8}.
  \end{itemize}
  \source{The number of values is the number for the fallback dataset.
    Your dataset gives a different number.}
\end{frame}
